% Transcription — fonds Alexandre Grothendieck, shelfmark 29, batch 8 (pages 141-160).
% Established from the facsimile archives/batches/29/batch-08.pdf (a mirror of
% https://grothendieck.umontpellier.fr/29.pdf), per the skill
% transcribe-grothendieck.
%
% Pass: Opus 5 (claude-opus-5), 2026-08-29 — first pass, unchecked against the
% pages by a human.
%
% Seventeen of the twenty pages are transcribed. Three are not:
%
%   - page 143, an otherwise blank sheet inscribed, in his hand, top right,
%     « Compléments SGA / 1960 ». It is the wrapper of the run that follows
%     and names it — page 148 opens « Compléments aux notes IHES » — so it
%     gets no \page{}, exactly as page 123 of batch 7 and page 111 of batch 6;
%   - pages 145 and 147, two sheets of a French typescript (a corollary
%     attributed to Hilbert on Jacobson rings, paginated "- 156 -"; a passage
%     on uniform structures and Cauchy filters, paginated "- 167 -"), on whose
%     backs the handwritten pages 146 and 148 are written. They carry nothing
%     in his hand, which is the criterion batches 5, 6 and 7 used for typed
%     versos: kept when annotated, skipped when not. Both are scanned upside
%     down for the same reason.
%
% Page 149 is kept although it carries almost nothing. Its verso is a letter of
% May 1961, but the recto side bears, in his hand, a struck word and the sketch
% X -> S with S' beside it — the base-change situation pages 148 and 150 are
% arguing about. A verso that continues the calculation is transcribed like any
% other page.
%
% The batch falls into seven runs.
%
%   141-142  the close of the fibre-product construction begun on pages 139-140
%            of batch 7 — the marked objects (P,xi), (Q,eta) and their base
%            change to S_1 — then z(G x G') = z(G) x z(G'), the pro-object
%            pi_1^C(S,xi), and two examples. Pencil, and the faintest writing
%            in the batch; the apparatus is correspondingly heavy here.
%   144-146  descent: the projections p_ij, p_i between pi', pi''_alpha,
%            pi'''_beta, the inverse images of E, and the cocycle relation
%            A_{alpha''} a''_beta A_{alpha'} = a'''_beta A_{alpha'''} a'_beta,
%            which page 144 boxes. Page 146 restates the same thing in terms of
%            paths and is then cancelled almost entirely.
%   148-150  « Compléments aux notes IHES »: the canonical homomorphism
%            pi_1(X x_S Y, c) -> pi_1(X,a) x_{pi_1(S,d)} pi_1(Y,b), its
%            surjectivity, the condition on S' making it bijective, the case of
%            a field, and the ladder of exact sequences that proves it.
%   151      the category of pi-sets, its fibred categories, and
%            H^i(F) = H^i(pi,E). The page is cancelled by three long diagonal
%            strokes.
%   152-156  « 3. Application »: Proposition 3.2 characterising X -> Y finite
%            étale as X = P/H for P galois, with the implication (i) => (iii
%            bis) credited on the page to J.-P. Serre; then Proposition 3.1 on
%            the existence of X/G and its étaleness over Y.
%   157-158  « Catégorie galoisienne » — axioms Gal 1 to Gal 6, in the numbering
%            he settles on after renumbering on the page — followed by the same
%            axiomatics stated for a category C with a full subcategory P and a
%            class of "non ramified" morphisms, 1° to 9°.
%   159-160  « Groupe fondamental dans les catégories »: a category with a class
%            (D) of morphisms of descent, principal bundles, structure groups.
%
% Page 152's title line is struck through twice over and defeated the pass; it
% is left as \ill{}. Pages 155 and 156 are pencil over a show-through and are
% the second hardest run after 141.

\documentclass[11pt,a4paper]{article}
% Le préambule commun à toutes les transcriptions du fonds Grothendieck.
%
% Il tient en une idée : **l'incertitude se compose, elle ne se résout pas.**
% Une transcription qui lisse un mot douteux a détruit la seule chose pour
% laquelle on la faisait — un lecteur doit pouvoir savoir, à chaque endroit,
% ce qui était sur le feuillet et ce qui a été ajouté. Les six macros qui
% suivent sont donc l'appareil critique, et non de la décoration.
%
% Les mêmes commandes sont reconnues par `scripts/render.mjs`, qui produit la
% vue de lecture du site. Une macro ajoutée ici doit l'être aussi là-bas,
% faute de quoi le rendu échouera bruyamment — ce qui est voulu : un rendu
% silencieusement faux est pire qu'un rendu absent.

\ProvidesPackage{grothendieck}[2026/08/08 Transcriptions du fonds Grothendieck]

\usepackage{amsmath,amssymb,amsthm}
\usepackage{mathrsfs}
\usepackage{stmaryrd}
% Les diagrammes commutatifs sont partout dans ce fonds : c'est la langue
% même de la géométrie algébrique de Grothendieck, et une transcription qui
% les rendrait en prose ne transcrirait rien.
\usepackage{tikz-cd}
% Français par défaut — c'est la langue des feuillets — mais l'anglais est
% chargé aussi : la traduction et le résumé sont des documents anglais, et la
% typographie française (espaces avant les deux-points) y serait une faute.
% Ces documents basculent par \selectlanguage{english} en tête de corps.
\usepackage[english,french]{babel}
\usepackage{fontspec}
\usepackage{geometry}
\geometry{a4paper,margin=25mm,marginparwidth=28mm}
\usepackage{marginnote}
\usepackage{xcolor}
% ulem, not soul. `soul`'s \st and \ul are built on a character-by-character
% scan that XeLaTeX's font machinery breaks: the document fails with an
% undefined control sequence rather than degrading. ulem does the same two jobs
% and is Unicode-engine safe. `normalem` stops it hijacking \emph, which the
% transcriptions use for Grothendieck's own emphasis.
\usepackage[normalem]{ulem}
\usepackage{hyperref}
\hypersetup{colorlinks=true,linkcolor=black,urlcolor=black,pdfborder={0 0 0}}

% --- Métadonnées du lot -----------------------------------------------------
% Reprises telles quelles par le script de rendu, qui les affiche en tête de
% la vue de lecture. Elles fixent ce que le fichier prétend couvrir : sans
% elles, une transcription flotte sans point d'attache dans un fonds de seize
% mille feuillets.

\newcommand{\folder}[1]{\gdef\grothFolder{#1}}
\newcommand{\batch}[1]{\gdef\grothBatch{#1}}
\newcommand{\leaves}[2]{\gdef\grothFirst{#1}\gdef\grothLast{#2}}
\newcommand{\foldertitle}[1]{\gdef\grothFoldertitle{#1}}
\gdef\grothFolder{?}\gdef\grothBatch{?}\gdef\grothFirst{?}\gdef\grothLast{?}\gdef\grothFoldertitle{}

% --- Le résumé --------------------------------------------------------------
%
% Il ouvre la lecture modernisée, sous le titre « Résumé », et il est composé
% en petit. Ce n'est pas un ornement : c'est le seul passage du document qui
% ne soit pas des mathématiques, et un lecteur doit pouvoir le reconnaître
% avant de l'avoir lu — puis le sauter s'il connaît déjà le sujet, ou s'y
% arrêter s'il ne le connaît pas. Le filet qui le suit marque la reprise du
% corps.

\newenvironment{resume}
  {\small\setlength{\parskip}{0.45em}}
  {\par\medskip\hrule\medskip}

% --- Mots-clés ---------------------------------------------------------------
%
% En anglais, en clôture du résumé de la lecture modernisée : le vocabulaire
% moderne sous lequel chercher ce que les feuillets construisent. Le manifeste
% du site (`npm run manifest`) les extrait pour étiqueter la cote entière —
% cette ligne est la source unique des tags, il n'y a pas de fichier à côté.
\newcommand{\keywords}[1]{\par\smallskip\noindent
  {\footnotesize\textsc{Keywords} --- \textit{#1}}\par}

% --- L'appareil critique ----------------------------------------------------

% Le numéro de feuillet, en marge. C'est l'ancre qui permet de régler un
% désaccord sur une formule en regardant *un* feuillet plutôt qu'une chemise
% de six cents. Le site s'en sert aussi pour tourner les pages du fac-similé
% quand on fait défiler la transcription.
\newcommand{\leaf}[1]{%
  \marginnote{\footnotesize\color{gray!70}\textsf{#1}}%
  \label{leaf:#1}\ignorespaces}

% Illisible : on ne devine pas. Un mot inventé qui se lit comme les autres est
% le pire résultat possible.
\newcommand{\ill}{\textcolor{red!60!black}{[\,\ldots\,]}}

% Lecture douteuse : le mot est donné, et signalé comme incertain.
\newcommand{\uncertain}[1]{\uline{#1}}

% Ajout de l'éditeur — une lettre manquante, un mot sous-entendu.
\newcommand{\add}[1]{\textcolor{blue!50!black}{[#1]}}

% Biffé par Grothendieck lui-même. Ce qu'il a rayé fait partie du document :
% une rature dit souvent où la pensée a changé de direction.
\newcommand{\struck}[1]{\sout{#1}}

% Note du transcripteur, sur l'état matériel du feuillet.
\newcommand{\note}[1]{\par\smallskip{\small\color{gray!60!black}\textit{[#1]}}\par\smallskip}

% Note marginale de Grothendieck, distincte du corps du texte.
\newcommand{\marginal}[1]{\par\smallskip\noindent\hspace{1em}%
  {\small\color{gray!60!black}#1}\par\smallskip}

% --- Titre ------------------------------------------------------------------

\AtBeginDocument{%
  \begin{center}
    {\large\bfseries Fonds Alexandre Grothendieck --- cote n\textsuperscript{o}~\grothFolder}\\[2pt]
    {\itshape\grothFoldertitle}\\[4pt]
    {\small feuillets \grothFirst--\grothLast\ (lot \grothBatch)}\\[2pt]
    {\footnotesize\color{gray!60!black}Transcription établie d'après le fac-similé numérisé
     par l'Université de Montpellier.\\
     Les lectures incertaines sont soulignées, les illisibles notées [\,\ldots\,],
     les ajouts entre crochets.}
  \end{center}
  \vspace{1em}}

\endinput


\folder{29}
\batch{8}
\pages{141}{160}
\dating{1959-1969}
\watermark{Édition de démonstration}
\foldertitle{Groupe fondamental [Autour de SGA 4 et SGA 7] : lettres (1959, 1967, 1969), tapuscrits et copies de tapuscrit annotés (s.d.), notes manuscrites (s.d.)}

\begin{document}

\section*{Fin de la construction du produit fibré ; $\pi_1^{\mathcal{C}}(S,\xi)$ (pages 141 à 142)}

\page{141}

\note{la page est au crayon et très pâle ; c'est la plus difficile du lot, et
l'appareil critique y est en conséquence plus lourd qu'ailleurs.}

\ldots{} la situation pour $G_1$, $G_1'$, $(P_1,\xi_1)$, $(Q_1,\eta_1)$ etc.
[\uncertain{obtenue}, \ill{} par le fait que \uncertain{$Q_1$} est vide,
$S_1$ \ill{} \uncertain{l'idée} \ill{}]. $P_1$ s'identifie alors
\add{à} $G_1'$ \marginal{(muni du marquage $\varepsilon_1'$)}, $Q_1$ à
\struck{$G_1'$} $G_1' \times_{S_1} G_1'$ muni de
$(\varepsilon_1', \varepsilon_1')$, \ill{}

$(P_1', \xi_1')$ s'identifie aussi à $(G_1', \varepsilon_1')$. Il suffit
de voir que $(P_1', \xi_1') \to (G_1',$ \ill{} $)$ est
\uncertain{l'inclusion diagonale} $G_1' \to G_1' \times_{S_1} G_1'$. On
\ill{}

\struck{\ill{}} \ldots{} \uncertain{relations}, et désignons par
$w$ l'isomorphisme $(P,\xi) \times_u G' \to (P,\xi) \times_v G'$,
$w_1$ l'isomorphisme \uncertain{déduit} par \uncertain{extension de la base}
à $S_1$, notre assertion signifie que $w_1$ transforme
\struck{\ill{}} \uncertain{la section} « naturelle » de
$P_1 \times_{v_1} G_1'$, i.e. \ill{}

\begin{tikzcd}[column sep=large]
P \arrow[r, "\mathrm{can}"] \arrow[rr, bend right=30, "\mathrm{can}"'] & P \times_u G' \arrow[r, "w"] & P \times_v G'
\end{tikzcd}

Or cela résulte de l'unicité de $P \to P \times_u G'$ \ill{}
\note{la dernière ligne, en bas de feuille et en partie biffée, n'a pas été
lue.}

\page{142}

\struck{\ill{}}

D'autre part on a
\[
\mathfrak{z}(G \times G') = \mathfrak{z}(G) \times \mathfrak{z}(G')
\]
et ainsi de \uncertain{même}. Les \uncertain{s-gpes} $G \in \mathcal{C}$ d'un
$G \in \mathcal{C}$ satisfaisant la condition minimale. \note{la phrase est
telle quelle sur la page : le même symbole $G$ y sert deux fois.}

\ill{} Ces propriétés formelles suffisent pour
construire un pro-objet de $\mathcal{C}$, au
moyen d'un système projectif d'épimorphismes, \uncertain{que nous désignerons}
par $\pi_1^{\mathcal{C}}(S, \xi)$.

\emph{Exemple I} $\mathcal{C} = $ \uncertain{un ensemble} \ill{} de groupes
définis par un groupe \ill{}. On retrouve le
groupe fondamental.

\emph{Exemple II} $S$ est un schéma \add{ou} \ill{}
d'un corps $k$. $\mathcal{C}$ provient d'une \ill{}

\section*{Données de descente : les conditions de cocycle (pages 144 à 146)}

\page{144}

\[
\begin{array}{ccc}
\pi' & \pi''_{\alpha} & \pi'''_{\beta} \\
S & S' & S'''
\end{array}
\]
\note{sous $\pi'''_{\beta}$ la page porte d'abord \struck{$S''$}, corrigé en $S'''$.}

\[
\begin{array}{ll}
p_{ij}^{\alpha\beta} : \pi'''_{\beta} \to \pi''_{\alpha} & \alpha = p_{ij}(\beta), \quad (ij) = \{(1,2), (2,3), (1,3)\} \\
p_i^{\alpha} : \pi''_{\alpha} \to \pi' & i = \{1, 2\}
\end{array}
\]

\struck{$p_i^{\alpha} p_{ij}^{\alpha\beta}$}

\[
\begin{array}{l}
p_1^{\alpha} p_{12}^{\alpha\beta} = \operatorname{ind}(a'_{\beta})\, p_1^{\alpha} p_{13}^{\alpha\beta} \\
p_2^{\alpha} p_{12}^{\alpha\beta} = \operatorname{ind}(a''_{\beta})\, p_1^{\alpha} p_{23}^{\alpha\beta} \\
p_2^{\alpha} p_{13}^{\alpha\beta} = \operatorname{ind}(a'''_{\beta})\, p_2 p_{23}
\end{array}
\]

\[
\begin{array}{ll}
p_{12}^{-1}(p_1^{-1}(E)) \to p_{12}^{-1}(p_2^{-1}(E)) & (p_1 p_{12})^{-1}(E) \to (p_2 p_{12})^{-1}(E) \\
p_{23}^{-1}(p_1^{-1}(E)) \to p_{23}^{-1}(p_2^{-1}(E)) & (p_1 p_{23})^{-1}(E) \to (p_2 p_{23})^{-1}(E)
\end{array}
\]

\begin{tikzcd}[column sep=small, row sep=small, nodes={font=\scriptsize}]
(p_1 p_{12})^{-1}(E) \arrow[r] \arrow[d] & (p_2 p_{12})^{-1}(E) = (p_1 p_{23})^{-1}(E) \arrow[d] \\
(p_1 p_{13})^{-1}(E) \arrow[r] & (p_2 p_{13})^{-1}(E) = (p_2 p_{23})^{-1}(E)
\end{tikzcd}

\note{les deux flèches du carré portent chacune une étiquette d'une seule
lettre, qui n'a pas été lue ; elles sont laissées nues plutôt que devinées.}

Pour tout $\beta$, on a \uncertain{ainsi}, en
posant $\alpha' = p_{12}(\beta)$, $\alpha'' = p_{23}(\beta)$,
$\alpha''' = p_{13}(\beta)$, \struck{\ill{}}

\begin{tikzcd}[column sep=large]
E \arrow[r, "A_{\alpha'}"] \arrow[d, "a'_{\beta}"'] & E \arrow[r, "a''_{\beta}", "\sim"'] & E \arrow[d, "A_{\alpha''}"] \\
E \arrow[r, "A_{\alpha'''}"'] & E \arrow[r, "a'''_{\beta}"', "\sim"] & E
\end{tikzcd}

\[
\boxed{\; A_{\alpha''}\, a''_{\beta}\, A_{\alpha'} = a'''_{\beta}\, A_{\alpha'''}\, a'_{\beta} \;}
\]

\page{146}

\note{page presque entièrement biffée : elle reprend en termes de chemins la
relation encadrée de la page 144, et l'abandonne. Ce qui subsit\uncertain{e}
de lisible est donné ici ; le reste est signalé, non deviné.}

Si on a choisi les \struck{$s''$ et $s'''$ tels que \ill{}}
\marginal{pour \uncertain{$s''' = \operatorname{diag}_{\theta_2}(s'')$}}
et \struck{\ill{}} \uncertain{$s''' = \operatorname{diag}_{g_3}(s')$}

\struck{\ill{}} \ill{} que par \ill{} $s' \in E$, désignant
par \ill{} $s'''$ \ill{}, $s'$ \ill{} de $E'$, \ill{} par l'application
diagonale, \struck{\ill{}} $s''$ \add{sont} les images de $s'$ et $s''$ par
\struck{les applications diagonales}, et \ill{}

\ldots{} chemins \struck{\ill{}} à des $s''$, $s'''$ \struck{\ill{}}
\marginal{\ill{} pour \ill{} $s'''$, $s'''$ que $s'' = s'$.}
chemins triviaux. (Alors il résulte de ce qui \ill{}
\[
g_{s''} = g_{s'''}\, g_{s'} \quad \text{i.e.} \quad g_{s''} = 1
\]
à la relation b), pour \ill{} la forme équivalente
b') \quad \uncertain{$g_{\bar{s}_i} = 1$}

\begin{tikzcd}[column sep=small, row sep=small]
& R_0 \arrow[d] & R_1 \arrow[d] & R_2 \arrow[d] \\
S & S' \arrow[l] & S'' \arrow[l] & S''' \arrow[l]
\end{tikzcd}

\note{le bas de la page ne porte que ce croquis, avec à côté les mots isolés
$\pi_1(S,\alpha)$, $\pi$ et $S$ ; il n'y a pas d'énoncé qui les relie.}

\section*{Compléments aux notes IHES : $\pi_1$ d'un produit fibré (pages 148 à 150)}

\page{148}

\emph{Compléments aux notes IHES}

\struck{Exercice}

\emph{Remarque.} La méthode de démonstration \ill{} en 1.7 \ill{}
\ill{} généralité \ill{} l'énoncé suivant : Soient $X$ et $Y$ sur
$S$, $X$ propre et \uncertain{séparable} \marginal{surjectif} \ill{},
et $Y$ \ill{}, \marginal{à \ill{} non vide.}
\struck{Désignons} un point géométrique $c$ de $X \times_S Y$, d'images
points géométriques $a$, $b$, $d$ dans $X$, $Y$, $S$, et un
homomorphisme canonique
\[
\pi_1(X \times_S Y, c) \to \pi_1(X, a) \times_{\pi_1(S, d)} \pi_1(Y, b).
\]
Cet homomorphisme est surjectif. \struck{Si \ill{}} \ldots{} de la bijectivité si
\struck{$Y$ a une section sur $S$, plus généralement si}
la condition suivante est vérifiée : Il existe
un $S'$ \struck{\ill{}} (non vide, \struck{\ill{}}) sur $S$, tel que
$Y \times_S S' = Y'$ ait une section sur $S'$ [i.e. \struck{il faut}
il suffit un $S$-morphisme $S' \to Y$] \struck{et \ill{}}
\struck{\ill{}} et que
$\pi_1(X') \to \pi_1(X)$ soit \uncertain{injectif}, \ill{}
$S'$ étale sur $S$ \struck{\ill{}}. Cet énoncé
s'applique en particulier \ill{} \ill{} où $S$
est le spectre d'un corps $k$, et donnera
\[
\pi_1(X \times_k Y) \simeq \pi_1(X) \times_{\pi_1(k)} \pi_1(Y)
\]
[$X$ propre universellement connexe sur $k$, $Y$ \struck{\ill{}}
\ill{} sur $k$ sans plus].

\page{149}

\struck{\ill{}}

\begin{tikzcd}[row sep=small, column sep=large]
X \arrow[d] & \\
S & S'
\end{tikzcd}

\note{seul ce croquis figure sur la feuille.}

\page{150}

\[
\begin{array}{ccc}
X & & Y \\
& S &
\end{array}
\]

$X$ propre séparable \struck{\ill{}} universellement connexe sur $S$,
$Y$ : tel qu'il existe \struck{\ill{}} un $S'$ \struck{\ill{}}
\ill{} non vide, fini, de \uncertain{surjectif} $S' \to S$, et une
section de $Y \times_S S'$ sur $S'$
[\uncertain{un peu} $=$ penser à un nouveau voisinage de $S$].

\begin{tikzcd}[column sep=small, row sep=small, nodes={font=\scriptsize}]
\pi_1(X \times_S Y) & & \\
\pi_1(X) \arrow[d, "\mathrm{surj}"'] & \pi_1(Y) \arrow[dl, "\mathrm{surj}"] & \pi_1(X') \arrow[d] \\
\pi_1(S) \arrow[rr] & & \pi_1(S')
\end{tikzcd}

\note{le trait entre $\pi_1(S)$ et $\pi_1(S')$ porte une pointe à chaque
extrémité sur la page.}

\begin{tikzcd}[column sep=small, row sep=small, nodes={font=\scriptsize}]
& \pi_1(X) \times_{\pi_1(S)} \pi_1(Y) \arrow[dl] \arrow[dr] & \\
\pi_1(X) \arrow[dr, "\mathrm{surj}"'] & & \pi_1(Y) \arrow[dl] \\
& \pi_1(S) &
\end{tikzcd}

\begin{tikzcd}[column sep=small, row sep=small, nodes={font=\scriptsize}]
& \operatorname{Im}(\pi_1(\bar{X}_s) \to \pi_1(X)) \arrow[d, no head] & & & \\
e \arrow[r] & \operatorname{Ker}[\pi_1(X) \to \pi_1(S)] \arrow[r, "\sim"] & \pi_1(X) \times_{\pi_1(S)} \pi_1(Y) \arrow[r] & \pi_1(Y) \arrow[r] & e \\
e \arrow[r] & \operatorname{Ker}[\pi_1(X \times_S Y) \to \pi_1(Y)] \arrow[u] \arrow[d, no head] \arrow[r] & \pi_1(X \times_S Y) \arrow[u] \arrow[r] & \pi_1(Y) \arrow[u] \arrow[r] & e \\
& \operatorname{Im}(\pi_1(\bar{X}_s) \to \pi_1(X \times_S Y)) & & &
\end{tikzcd}

\note{les traits verticaux sans pointe rendent les deux signes $\parallel$ de
la page.}

\begin{tikzcd}[column sep=small, row sep=small]
\bar{X}_s \arrow[r] \arrow[d, no head] & X \times_S Y \arrow[d] \\
\bar{X}_s \arrow[r] & X \arrow[d] \\
& S
\end{tikzcd}

\begin{tikzcd}[column sep=small, row sep=small]
\pi_1(\bar{X}_s) \arrow[r] \arrow[d, no head] & \pi_1(X \times_S Y) \arrow[d] \\
\pi_1(\bar{X}_s) \arrow[r] & \pi_1(X) \arrow[d] \\
& \pi_1(S)
\end{tikzcd}

\section*{Le topos des $\pi$-ensembles et la cohomologie de $\pi$ (page 151)}

\page{151}

\note{la page entière est barrée de trois longs traits diagonaux.}

$\pi$ un groupe.

\begin{enumerate}
\item[1)] $\mathcal{E}^{\pi}$ catégorie des ensembles \ill{} où $\pi$ opère
\marginal{à gauche}, \struck{\ill{}} \ill{} par les
\uncertain{morphismes surjectifs}, objet final $S = e$, objet
\struck{\ill{}} initial $\emptyset$, objets minimaux \uncertain{ceux qui}
$\neq \emptyset$
\[
\pi_s = \pi, \qquad \operatorname{Aut}(T) = \operatorname{Aut}(\pi/s) \simeq \pi
\quad \text{par } x \mapsto x g^{-1}
\]
\marginal{i.e. \ill{} que $S \to e$ soit \uncertain{surjectif} \ill{}}
\end{enumerate}

Faisceau d'ensembles sur $\mathcal{E}$ $\simeq$ \marginal{catégorie des}
ensembles finis où $F$ \uncertain{opère} à gauche
\[
E = F(\pi_s), \qquad F(U) = \operatorname{Hom}_{\pi}(U, E)
\]
\note{un membre biffé, non lu, précède $\operatorname{Hom}_{\pi}(U,E)$ : \struck{\ill{}}}
\marginal{où on fait opérer $\pi$ via les opérations à gauche sur $\pi_s$.}

$\mathcal{F}^{\pi}$ catégorie fibrée en \ill{} sur les \ill{} dans
$\mathcal{E}_{\pi}$
\[
\mathcal{F}^{\pi}_e \simeq \mathcal{F}^{\pi}_s \simeq \mathcal{E}^{\pi}_{\mathrm{Ens}}
\]
\[
\mathcal{F}^{\pi}_T \simeq (\mathrm{Ens}) \qquad \text{car} \qquad
\mathcal{E}_{/T} \simeq (\mathrm{Ens}) = \mathcal{E}^{\circ}
\]
\marginal{foncteur fibre $e \in \pi_s = T$}
\marginal{changement de base \ill{}, se déduit \ill{} par les opérations de $\pi$}

Plus précisément, si
\[
U = \coprod_i \pi/\pi_i \qquad \text{on a} \qquad \mathcal{F}_U \simeq \prod_i E^{\pi_i}
\]
\marginal{décomposition en \uncertain{indécomposables}}

\[
\begin{array}{lll}
H^0(F) = \Gamma(F) = F(e) & \text{correspond à} & E^{\pi} = \operatorname{Hom}(\pi_s, E) \\
H^1(F) & \simeq & H^1(\pi, E) \\
H^i(F) & \simeq & H^i(\pi, E)
\end{array}
\]
la dernière ligne vaut si $F$, $E$ sont \uncertain{commutatifs}.

\section*{Application : caractérisation des morphismes finis étales (pages 152 à 156)}

\page{152}

\begin{enumerate}
\item[3.] \uncertain{Application}. La caractérisation des \struck{\ill{}}
\marginal{variétés} de la \ill{} \struck{\ill{}} \ill{}
\note{le titre est biffé et resurchargé deux fois ; il n'a pas été lu.}
\end{enumerate}

\marginal{\ill{}, $G_n$}

\emph{Prop 3.2.} Soit $f : X \to Y$ un morphisme fini, \ill{}

Conditions équivalentes
\begin{enumerate}
\item[(i)] $f$ est étale, et le rang de $X$ \ill{} \ill{} sur $Y$ \ill{}
localement égal à $n$.
\item[(ii)] Il existe un $Y_1 \to Y$ fidèlement plat tel que
$X \times_Y Y_1$ soit $Y_1$-isomorphe à $Y_1 \times I_n$.
\item[(ii bis)] Kif -- kif, avec $Y_1$ \struck{\ill{}} un revêtement
galoisien (de groupe $G_n$ \uncertain{il me semble}).
\item[(iii)] $X$ est isomorphe à un $P$\struck{\ill{}}$/H$, où $P$ est un
revêtement étale galoisien de $Y$, de groupe $G$, et $H$ un sous-groupe de $G$
\marginal{\ill{} d'indice \ill{}}
\item[(iii bis)] Kif kif, avec $G = G_n$, $H = G_{n-1}$
(stabilisateur d'un point de $I_n$).
\end{enumerate}

\marginal{\struck{(ii) $\Rightarrow$ (iii bis)} \struck{\ill{} de \ill{}}
\struck{(iii) $\Rightarrow$ (iii bis) \uncertain{immédiat}} \struck{\ill{}}}

(i) $\Rightarrow$ (iii bis) (la démonstration \struck{qui suit}
\marginal{\ill{} que la suite} due à J.-P. Serre). Le fait \ill{}
$G_n$ \ill{} $X \times_Y \cdots \times_Y X$ \ill{} par \uncertain{permutation}
des facteurs, \ill{} par $f$, \ill{} évidente. L'assertion des points
d'inertie est \struck{fausse} \marginal{viennent de composantes} connexes
(car si $s \in G_n$ est dans le \ill{}
d'inertie d'un point $\alpha$, il opère trivialement \ill{}

\page{153}

\ldots{} \ill{} la composante connexe \ill{} est réduit \add{au point}). Soit
$P$ son complémentaire. Il est stable \marginal{un ouvert et fermé de $Y$}
\ill{} $G_n$, qui y opère \ill{}. De plus, on doit
\ill{} \uncertain{pour tout} $\gamma \in \mathcal{C}$, \struck{l'étale}
\marginal{la fibre} $P$ sur $Y$ \ill{} \ill{}

\ldots{} \uncertain{permutation} \ill{} triviale de $I_n$ : il \ill{}
les \uncertain{points d'inertie} \struck{\ill{}} distincts de $E$,
\marginal{il y a $n!$ \ill{}}
En vertu de \emph{prop. 2.5}, $P \to Y$ \ill{} un
revêtement étale galoisien de $Y$ de groupe $G_n$.

D'autre part, \struck{\ill{}} la \struck{\ill{}} \marginal{dernière} \ill{}
\uncertain{$pr_{12}$}$\, : X^n \to X$ \struck{étale} est invariante sous
le stabilisateur $H = G_{n-1}$ de $\{n\}$, d'où
\struck{\ill{}} \ill{}

\page{154}

\ldots{} revenir au cas où $G$ opère transitivement sur
l'un \ill{} des composantes connexes. Soit $X_0$ \ill{}
d'elles, $G_0$ son stabilisateur, alors il est
évident que $X$ s'identifie à $X \times_{G_0} G$,
et $X/G \simeq X_0/G_0$.

\emph{Enfin}, \struck{on \ill{} peut supposer $X_0$ connexe. Soit $H$}
\struck{\ill{} soit $G_i$ le groupe d'inertie, il \ill{}}
\struck{qui opère trivialement \ill{}}
Donc on \ill{} que $G_0$ opère \ill{} fidèlement. Mais alors il
opère \emph{sans inertie}
\[
X \to X/G
\]
\ill{}
donc reste : voir que $X/G$ est \uncertain{étale} sur $Y$.

(Faisant une extension fidèlement \ill{} de la base $Y$ \ill{},
on est ramené au cas où l'extension résiduelle $k(x)/k(y)$ est
\uncertain{triviale}. Cela \ill{}, on voit que le diagramme
$\hat{\mathcal{O}}_y \to \hat{\mathcal{O}}_{x'}$ \ill{})

\begin{tikzcd}[column sep=small, row sep=small]
X \arrow[dr] & \\
& X' = X/G \arrow[d] \\
& Y \ni y
\end{tikzcd}

\begin{tikzcd}[row sep=small]
k(x') \arrow[d, no head] \\
k(x) \arrow[d, no head] \\
k(y)
\end{tikzcd}

\note{la tour est dessinée dans cet ordre sur la page ; les inclusions vont en
sens inverse, $k(y) \subset k(x') \subset k(x)$.}

\[
\hat{\mathcal{O}}_y \to \hat{\mathcal{O}}_{x'} \xrightarrow{\ \sim\ } \hat{\mathcal{O}}_x
\]

\emph{Corollaire} Soit $X$ étale sur $Y$ en $x$, $g : X \to X'$ un
$Y$-morphisme étale en $x$, alors $X'$ est
étale sur $Y$ en $g(x)$.

\page{155}

\ldots{} \uncertain{car} il est surjectif et radiciel, donc \ill{} pour tout
$\gamma \in Y$, la fibre géométrique de $P/H$ \struck{\ill{}} \ill{}
d'avoir \uncertain{bijectivement} \ill{} de $X$ en $\gamma$. Ce qui
\ill{} un point fini, d'après l'application
précédente.

\uncertain{On a ainsi} (iii bis) $\Rightarrow$ (iii), et (iii) $\Rightarrow$ (i)
d'après un \emph{lemme} précédent, déjà \uncertain{établi}.
Puis l'équivalence de (i) (iii bis) (iii).
De plus (iii) \struck{\ill{}} $\Rightarrow$ (ii bis), comme on \ill{}
en prenant $Y_1 = P$, et utilisant la formule
\[
(X \times G)/H \simeq X \times G/H,
\]
et (ii bis) $\Rightarrow$ (i). Enfin,
(ii bis) $\Rightarrow$ (i) est immédiat [car la condition
\struck{\ill{}} \ill{} si $X$ est étale par un \ill{}].

\emph{Cor.} \ill{} \ill{} $A$ sur $Y$, $X_1$ \ill{}
\ill{} $A_1$, si $A_1$ est \uncertain{libre}, il en est
\ill{} $= $ de $A$ comme \ill{}. Si
les fibres de $X_1$ sur $Y_1$ sont
\marginal{(distinctes de celles de $X$ sur $Y$ car \ill{})} \ill{}
\ill{} celles de $X$ sur $Y$. D'où \ill{}

\page{156}

\emph{Proposition 3.1} (\ill{}) $X$ étale sur $Y$ \struck{\ill{}}
\uncertain{quasi-compact}, $G$ un
groupe fini \marginal{opérant} \struck{dans} $Y$-automorphismes \struck{\ill{}}.
\begin{enumerate}
\item[(i)] \struck{si} chacun des points de $X$ \ill{} $G$ \ill{}
\struck{\ill{}} \ill{} \ill{} et pour \ill{}
\item[(ii)] \struck{\ill{}} $X/G$ existe, et $X \to X/G$ \ill{}
et $X/G \to Y$ \ill{} \emph{étale}.
\end{enumerate}

Si $s \in G$ \struck{\ill{}} \ill{} le groupe
d'inertie \ill{} \ill{} point $x \in X$, il \struck{laisse}
\marginal{qui laisse un point} \struck{fixe} \ill{}
(ii), \ill{} l'un $X_0$ \ill{}

\ldots{} des $x \in X$ tels que $s \in G_i(x)$, \ill{}
\ill{} \ill{}, \ill{} fini. \struck{\ill{}} \ill{} l'un des points \ill{}
i.e. $G_i(x) \neq e$ \ill{} la réunion des $X(s)$, $(s \neq e)$,

(i) \ill{} évident. \ill{} \ill{}, $X$ étant
\uncertain{quasi-affine} sur $Y$, $X/G$ existe [\ill{}]
\struck{\ill{}} \ill{} en faisant \ill{} ce que \ill{}.
Soit $I$ l'ensemble des composantes connexes de $X$.
$G$ opère sur $I$, \struck{et $X$ est réunion des}
\ldots{} $X_{(\alpha)}$, \struck{$\alpha \in$} $I/G$, obtenus en
posant
\[
X_{(\alpha)} = \bigcup_{i \in \alpha} X_i, \quad \text{stables sous } G.
\]
donc
\[
X/G = \bigcup X_{(\alpha)}/G
\]
\ill{}

\section*{Catégories galoisiennes (pages 157 à 158)}

\page{157}

\emph{Catégorie galoisienne}
\begin{enumerate}
\item[(i)] Une catégorie $\mathcal{C}$
\item[(ii)] Un foncteur $F : \mathcal{C} \to$ Ensembles finis
\end{enumerate}

\marginal{Pour \ill{} \ill{} axiomes l'existence d'un $F$ \ill{} :
\struck{unicité} \ill{} $F$ permet \ill{} finis \ill{}
$F(X) \to F(Y)$ \ill{}, $X$ isom. $\Rightarrow X \to Y$.}

Axiomes :

\marginal{aux sommes, aux produits, en \ill{} quelconques finis, et \ill{}
les qui \ill{}}

\begin{enumerate}
\item[(Gal 1)] Dans $\mathcal{C}$, les sommes et les produits
\struck{de deux objets} existent \marginal{[d'où l'existence de $0$ et $1$]}

\item[\struck{(Gal 3)}] \struck{$X \in \mathcal{C}$, $F(X) = \emptyset
\Rightarrow X$ est une unité à gauche (soit $0$)}

\item[(Gal 4)] $F$ permute à la somme et au produit
\end{enumerate}

$X \in \mathcal{C}$ se dit \struck{indécomposable} \marginal{connexe} si
$X = X' \sqcup X''$ implique $X' \simeq 0$ ou $X'' \simeq 0$.

\begin{enumerate}
\item[(Gal 2)] $X \in \mathcal{C}$ \struck{\ill{}} implique
$X \simeq \coprod_{i \in I} X_i$ \marginal{$X_i \neq 0$}, et
\struck{la} \ill{} décomposition $X \simeq X' \sqcup X''$ implique
$X' \simeq \coprod_{i \in I'} X_i$, $X'' = \coprod_{i \in I''} X_i$,
$I' \cap I'' = \emptyset$, $I' \cup I'' = I$. Les $X_i$, déterminés de
façon essentiellement unique (avec leur $X_i \to X$)
sont appelés les \struck{composants indécomposables} \marginal{connexes} de
$X'$ \marginal{[N.B. $X = 0 \Longleftrightarrow I = \emptyset$]}

\item[(Gal 3)] Soit $f : X \to Y$ avec $Y$ \struck{indécomposable}
\marginal{connexe $\neq 0$}. Alors les composantes \struck{$X$} $X_i$ de $X$
telles que $f$ induise un \emph{isom.} $f_i : X_i \xrightarrow{\sim} Y$ sont en
correspondance biunivoque avec les \emph{sections} $g : Y \to X$
\struck{de $f$} [i.e. telles que $fg = \mathrm{id}_Y$], $g$ étant la composée
\struck{\ill{}}
\[
Y \xrightarrow{\ f_i^{-1}\ } X_i \xrightarrow{\ \text{inj. can}\ } X
\]
[en d'autres termes, \struck{toute} \ill{} section se factorise]
\marginal{[N.B. il suffit de la vérifier si $X = Z \times Y$, $f = p_2$]}
\marginal{$X$, $Y$ quelconques}

\item[(Gal 5)] Sous les conditions de Gal 3, si \struck{\ill{}} $f : X \to Y$
induit $F(X) \xrightarrow{\sim} F(Y)$, c'est un isomorphisme
[implique $F(X) = \emptyset \to X = 0$]
\marginal{[il suffit de vérifier Gal 5 quand $X$ est une composante connexe
d'un produit $Z \times Y$, $f$ induit par $p_2$]}

\item[(Gal 6)] Soit $X \in \mathcal{C}$, $X$ connexe $\neq 0$, et galoisien,
\struck{i.e.} \struck{Alors} $G = \operatorname{Aut}(X)$ opère
\emph{transitivement} dans $F(X)$. \struck{\ill{}} Alors pour tout
\struck{sous-groupe} $H$ de $G$, $X/H$ dans $\mathcal{C}$ existe, et
\[
F(X/H) \simeq F(X)/H.
\]
\end{enumerate}

\page{158}

Catégorie $\mathcal{C}$, sous-catégorie $\mathcal{P}$ (pleine), morphismes $f$
« non ramifiés »

\begin{enumerate}
\item[1\textsuperscript{o}] Dans $\mathcal{C}$, les sommes finies existent,
ainsi que les produits fibrés.

\item[4\textsuperscript{o}] Dans $\mathcal{C}$, le \marginal{composé} produit de
deux morphismes non ramifiés \marginal{(des morphismes identiques)} où l'image
inverse d'un morphisme non ramifié sont non ramifiés, \ill{}

\item[5\textsuperscript{o}] \struck{\ill{}} $f : X \to Y$ \struck{est} un
\ill{} ramifié $\Longleftrightarrow$ \ill{} ($X$ : les composantes) \ill{}
\struck{les sommes} \ill{}

\item[2\textsuperscript{o}] Dans $\mathcal{C}$, il y a des composantes connexes.

\item[3\textsuperscript{o}] Les points $\xi \in \mathcal{P}$ sont connexes
(donc $\operatorname{Hom}(\xi, X \sqcup Y) = \operatorname{Hom}(\xi, X) \sqcup
\operatorname{Hom}(\xi, Y)$)

\item[6\textsuperscript{o}] Pour tout $f : X \to Y$ non ramifié,
\marginal{$Y$ connexe,} et $\xi \in \mathcal{P}$,
$\operatorname{Hom}(\xi, X) \to \operatorname{Hom}(\xi, Y)$ \ill{}

\item[7\textsuperscript{o}] Si \struck{$X$} et $Y$ sont connexes, et
\uncertain{les fibres} de $\operatorname{Hom}(\xi, X) \to
\operatorname{Hom}(\xi, Y)$ \marginal{surjectif} \ill{}
réduites à un élément, alors $f$ est un \emph{isom.}
(on suppose $f$ non ramifié, $\xi \in \mathcal{P}$)

\item[8\textsuperscript{o}] Si $f : X \to Y$ non ramifié, $Y$ connexe, alors
\struck{pour toute} \marginal{les} section $s : Y \to X$ est un \emph{isom.} de
$Y$ sur une composante connexe de $X$.

\item[9\textsuperscript{o}] Soit \struck{$G$} \marginal{\uncertain{groupe fini}}
\struck{opérant dans} $X$ \ill{} \ill{} $f : X \to Y$ non ramifié,
\struck{\ill{}} $X$ connexe, $\xi \in \mathcal{P}$,
$e \in \operatorname{Hom}(\xi, Y) = Y^{\xi}$ ; \ill{}
$G = \operatorname{Aut}_Y(X)$ opère \emph{transitivement} dans
$(f^{\xi})^{-1}(e)$. Alors pour tout
\uncertain{sous-groupe} $H$ de $G$, $X/H$ \ill{} \marginal{dans $\mathcal{C}$}
\ill{} un revêtement non ramifié de $Y$. \struck{\ill{}}
\end{enumerate}

\emph{N.B.} deux \uncertain{éléments} de $\operatorname{Hom}(\xi, X)$ \ill{}
\ill{} même \ill{} image par $X \to X/H$ sont \ill{} mod $H$.

Sous ces conditions, pour $x \in \operatorname{Hom}(\xi, X)$ on définit
$\pi_1(X, x)$. Je ne changerai pas \struck{\ill{}} \ill{} \ill{} $\xi$
\struck{\ill{}} \ill{} $\xi'$ \ill{} $\xi' \to \xi$ soit donné.
(c'est un foncteur \emph{covariant} en $(X, x)$)

\begin{tikzcd}
\xi \arrow[r, "x"] \arrow[d] & X \arrow[d] \\
\xi' \arrow[r, "x'"'] & X'
\end{tikzcd}

définissant $\pi_1(X, x) \to \pi_1(X', x')$. \emph{Etc.}

\section*{Groupe fondamental dans les catégories (pages 159 à 160)}

\page{159}

\emph{Groupe fondamental dans les catégories.}

$\mathcal{C}$ une catégorie, avec une famille $(D)$ de morphismes
dans $\mathcal{C}$, satisfaisant les conditions usuelles
\begin{enumerate}
\item[(i)] \uncertain{Contiennent l'identité} et \uncertain{stables} par
composition
\item[(ii)] \ill{} que dans \ill{} changement de base \marginal{$P$}
\ill{} les \ill{} de \ill{}
\item[(iii)] \struck{\ill{}} \ill{} sont des morphismes de descente
[\ill{} des \ill{} \emph{universels stricts} si \ill{} \ill{} et \ill{}
\ill{} produits existent].
\end{enumerate}

Enfin, on suppose
\begin{enumerate}
\item[(iv)] les \uncertain{morphismes} qui « descendent » par $g \in (D)$,
pour un $f \in (D)$ [mais on ne \uncertain{suppose} pas qu'on puisse faire
descendre tout morphisme par un $f \in (D)$].
\end{enumerate}

\marginal{(v) La projection $X \times I \to X$ est $\in (D)$.}

Dans la suite, $D$ est supposé donné une fois pour toutes.
\ill{} dans $\mathcal{C}$ les sommes finies \marginal{et \ill{} quelconques}
existent \ill{}

\emph{Exemple} Catégorie $\mathcal{C}$ des préschémas. \struck{La catégorie}
Comme $(D)$ les morphismes \emph{plats} et \emph{propres}.
[ou \emph{plats} et \emph{quasi-propres}].

On va développer une théorie particulière de
fibrés principaux [\uncertain{utilisant} les conditions générales
\struck{\ill{}} d'une théorie \ill{} fibrés principaux, cf. \ill{}].

\emph{Groupes structuraux :} $X \times G$, $G$ un groupe \uncertain{abstrait}
[ici fini \uncertain{quelconque}]

\emph{Fibrés principaux sur $X$, de groupe $G$ :} les fibrés
\uncertain{formés} ainsi, avec leur $P \to X \in (D)$.

\page{160}

\emph{Espace : opérateurs admissibles :} Les morphismes $E \to X$ sont
$\in D'$.

\emph{Hom admissibles de groupes structuraux :} \ill{} dérivé
de \ill{} \ill{} d'un \ill{} de \ill{}

\emph{Hom Inclusion admissible :} les \uncertain{morphismes} \ill{}
de groupes ordinaires.

\end{document}
